%% The guideline calls this chapter "Theoretical Approach"; rename it to suit your work.
\chapter{Methodology}
\label{ch:methodology}

Describe the theory, design and methods of your work in enough detail for someone else to
repeat it.

\section{Equations}
Write equations with LaTeX, never paste them as images (Guideline 6b). Each equation sits
on its own line and is numbered within the chapter:
\begin{equation}
  \mathrm{nDCG}@k = \frac{1}{\mathrm{IDCG}@k}\sum_{i=1}^{k}\frac{2^{rel_i}-1}{\log_2(i+1)}
  \label{eq:ndcg}
\end{equation}
Refer to it in the text as equation~\ref{eq:ndcg}.

\section{Figures}
Figures are numbered within the chapter and their captions are centred underneath, in
10~pt type (Guideline 6a). Figure~\ref{fig:pipeline} is drawn with TikZ; for an image file,
put it in the \texttt{figures} folder and use \verb|\includegraphics|.

\begin{figure}[htbp]
  \centering
  \begin{tikzpicture}[node distance=1.2cm,
      box/.style={draw,rounded corners,minimum width=2.8cm,minimum height=1cm,
                  align=center,font=\small}]
    \node[box] (data) {Data};
    \node[box,right=of data] (method) {Method};
    \node[box,right=of method] (eval) {Evaluation};
    \draw[->,thick] (data) -- (method);
    \draw[->,thick] (method) -- (eval);
  \end{tikzpicture}
  \caption{The three stages of the study}
  \label{fig:pipeline}
\end{figure}

\subsection{Sub-sections}
Sub-sections are numbered 3.2.1, 3.2.2 and so on. The guidelines allow no deeper
numbering; use bullet points instead.

\section{Summary}
Summarise the method.
